\chapter{Dense Reward and Credit Assignment}
\label{sec:reward}

\section{Syntactic Reward}
\label{subsec:syntactic-reward}

For an action \(a\) generated at state \(s\), let \(L_{\text{orig}}(a)\) be the list of clean proof lines in
the original generated action, excluding blank lines and comment-only lines. Let \(L_{\text{patch}}(a)\)
be the clean proof lines after Sorrifier recovery. The system compares these two lists using sequence
matching. A line is counted as surviving only if it appears unchanged and in the same relative order in
the patched proof.

Let \(N_{\text{orig}} = |L_{\text{orig}}(a)|\), and let \(N_{\text{surv}}\) be the number of surviving lines.
After the patched script is checked by Lean, the system also counts how many surviving lines correspond
to unused or ineffective code. We denote this number by \(N_{\text{dead}}\). These are lines that survived
the recovery process but do not appear to contribute useful executable proof structure, for example
because Lean reports that a declaration is unused or that a tactic has no effect. By definition,
\(0 \leq N_{\text{dead}} \leq N_{\text{surv}}\).

The syntactic reward is:
\[
r_{\mathrm{env}}(s,a) =
\begin{cases}
\dfrac{N_{\text{surv}} - N_{\text{dead}}}{N_{\text{orig}}},
& N_{\text{orig}} > 0, \\
0, & N_{\text{orig}} = 0.
\end{cases}
\]

This reward measures how much of the generated Lean code remains usable after recovery. Lines removed
or replaced by placeholders do not receive credit. Surviving lines that are later identified as unused or
ineffective are also discounted. Therefore, \(r_{\mathrm{env}}\) is a syntactic validity signal, not a proof
correctness signal. A high value means that much of the generated code survived Lean checking after patching; 
it does not mean that the action solves the proof state.

\section{Dependency-aware Structural Reward}
\label{subsec:dep-reward}

Syntactic validity alone is not enough to judge the quality of a generated action. Both skeleton actions
and local proof actions may introduce local declarations through constructs such as \texttt{have} or
\texttt{let}. In a skeleton action, these declarations often correspond to child subgoals that are solved by
child search branches and stitched back into the parent scaffold. In a local proof action, they are proof
steps inside a single generated proof body.

In either case, the model may introduce declarations that compile but do not help the final proof. It may
also leave declarations containing placeholders, which makes the resulting proof incomplete. To estimate
whether generated declarations actually contribute to the proof that was found, GammaZero performs
dependency analysis over the elaborated Lean proof expression.

This signal is more semantic than the syntactic reward in Section~\ref{subsec:syntactic-reward}: it is computed
from Lean's elaborated expression tree rather than from surviving source lines. It therefore asks whether
a declaration is used by the checked proof term, not merely whether the text around that declaration
compiled. It is still not a theorem about mathematical necessity or proof minimality; it is an attribution
signal for the concrete proof found by the search.

A separate expression-dumping process exports the Lean 4 expression tree into JSON. The dependency
analyzer traverses nodes such as \texttt{bvar}, \texttt{fvar}, \texttt{app}, \texttt{lam}, \texttt{forallE}, and
\texttt{letE} to determine whether each generated declaration is referenced by the final proof expression.
The traversal is De Bruijn-aware: when the analyzer enters a binder such as \texttt{lam}, \texttt{forallE},
or \texttt{letE}, it shifts the target De Bruijn index so that bound-variable references are attributed to
the correct declaration. The analyzer also detects \texttt{sorryAx} in the expression tree.

Combining usage detection with placeholder detection gives four structural cases:
\begin{itemize}
    \item \textbf{Used and solved:} the declaration has a solved proof body and is used by the elaborated proof. This is the core useful case.
    \item \textbf{Unresolved and used:} the declaration still contains a placeholder and is used by the proof, so the proof is incomplete.
    \item \textbf{Solved but unused:} the declaration is complete but redundant for the selected proof.
    \item \textbf{Unresolved and unused:} the declaration is unfinished and unused by the selected proof.
\end{itemize}

Let \(n_{\mathrm{used,solved}}\), \(n_{\mathrm{unused,solved}}\), and
\(n_{\mathrm{unused,unresolved}}\) denote the numbers of used-and-solved, solved-but-unused, and
unresolved-and-unused declarations. The general dependency
reward is:
\[
r_{\mathrm{dep}}(s,a) =
\frac{n_{\mathrm{used,solved}}}
{n_{\mathrm{used,solved}}
 + w_{\mathrm{unused,solved}} n_{\mathrm{unused,solved}}
 + w_{\mathrm{unused,unresolved}} n_{\mathrm{unused,unresolved}}},
\]
where \(w_{\mathrm{unused,solved}} \ge 0\) penalizes solved but unused work and
\(w_{\mathrm{unused,unresolved}} > w_{\mathrm{unused,solved}}\) penalizes unresolved unused work more
strongly. If the denominator is zero and the action introduces no local declarations, the reward defaults
to \(1.0\). Unresolved-and-used declarations are treated as
structural failures because the final proof still depends on a placeholder-containing declaration.

The formula is only a heuristic for ranking search actions. It favors used-and-solved declarations and penalizes unused or unfinished ones.
The numerator counts declarations that are both used and solved. This conjunction is essential: a declaration is not useful for
the final proof merely because it is solved, and it is not acceptable merely because it is used. It must be
both. The denominator counts those declarations plus weighted
forms of waste. A solved-but-unused declaration is mildly harmful because it consumes search effort but
does not block final verification. An unresolved-and-unused declaration is worse because it leaves
unfinished code in a placeholder-compilable script. The dependency score therefore helps distinguish
compact useful structure from redundant or bloated proof scripts, while final correctness still depends
only on Lean verification without placeholders.

This analysis is tied to the proof that GammaZero actually found. If a declaration is unused in the
selected stitched proof, that does not mean it is mathematically useless in every possible proof. It only
means that this search trajectory did not need it.

\section{AND/OR Backup}
\label{subsec:and-or-backup}

After search stops or the root is closed, GammaZero backs up values through the explored AND/OR graph.
States are OR nodes: a state can be solved by choosing one successful outgoing action. Actions are AND
nodes: a skeleton action succeeds only when all of its child proof states are solved.

The implemented backup uses a single action-value rule for both local proof actions and skeleton
actions. Solved actions are credited through the dependency reward and, for skeletons, through solved
child values.

For an action \(a\) applied to state \(s\), let \(\mathcal{C}(a)\) be the set of child states created by the
action. For a local proof action, \(\mathcal{C}(a)=\emptyset\). We define the empty-child future value to
be zero:
\[
\min_{s' \in \emptyset} V(s') = 0.
\]
The action value is:
\[
Q(s,a) =
r_{\mathrm{env}}(s,a)
+
r_{\mathrm{dep}}(s,a)
+
\mathds{1}_{\mathrm{solved}}(s,a)\,
\gamma
\min_{s' \in \mathcal{C}(a)} V(s').
\]

Here, \(r_{\mathrm{env}}\) is the syntactic preservation reward from the Sorrifier, and
\(r_{\mathrm{dep}}\) is the dependency reward from the elaborated proof expression. The indicator
\(\mathds{1}_{\mathrm{solved}}(s,a)\) is \(1\) only when the action solves its required proof state. For a
local proof action, this means the current target is closed without \texttt{sorry} or \texttt{admit}. For a
skeleton action, this means every child state in \(\mathcal{C}(a)\) is solved.

This rule has two useful special cases. For a local proof action, the child set is empty, so the future term
drops out:
\[
Q(s,a)=r_{\mathrm{env}}(s,a)+r_{\mathrm{dep}}(s,a).
\]
For a skeleton action that is not fully solved, the solved indicator is zero, so the action keeps only its
local credit. In practice, \(r_{\mathrm{dep}}\) is zero when the skeleton has not produced a solved
dependency structure. When the skeleton is solved, it also receives the discounted value of its weakest
child branch:
\[
Q(s,a)=
r_{\mathrm{env}}(s,a)
+
r_{\mathrm{dep}}(s,a)
+
\gamma \min_{s' \in \mathcal{C}(a)} V(s').
\]

The \(\min\) operator implements the AND part of the graph: a decomposition is limited by its weakest
child state. This prevents the backup from giving high delayed value to a skeleton that solves only easy
subgoals while leaving an essential one open.

For an OR node \(s\), the state value is the best available outgoing action:
\[
V(s)=\max_{a \in \mathcal{A}(s)} Q(s,a),
\]
where \(\mathcal{A}(s)\) is the set of actions expanded from \(s\). The graph therefore uses max backup
over alternative actions and min backup over the conjunctive children of a skeleton action.

\section{Reward Computation and Training Data Generation}

The implemented system produces scored proof-search trajectories rather than an end-to-end optimized
policy checkpoint. Each sampled action is executed, optionally recovered, inserted into the AND/OR graph,
and assigned syntactic and delayed structural rewards according to the formulations above. The resulting records
contain the serialized proof state, action type, generated Lean code, recovery result, child proof states,
graph status, priority-queue metadata, skeleton commitment metadata, and final action value.

These records can be used as training data for later policy optimization. Local proof actions and
skeleton actions should be grouped separately because their value is observed through different graph
structures: local proof actions have no child value term, whereas skeleton actions receive delayed child
value only when all generated child subgoals are solved. Therefore, comparing local proof actions and
skeleton actions in the same optimization group would mix different reward scales.

\section{Trajectory Export and Final Verification}

For a successful root theorem, GammaZero reconstructs the proof by recursively inserting solved child
proof bodies into the child subgoals of the committed parent skeletons. Reserved skeletons and failed
actions remain in the search record, but they are not part of the selected final proof. After insertion,
dependency analysis may remove solved declarations that are unused by the elaborated proof expression.
The resulting theorem is submitted to Lean again. The theorem is accepted as solved only if this final
verification succeeds without real \texttt{sorry} or \texttt{admit}.

The same trajectory records can later support policy optimization, but this thesis stops at verified
search, reward computation, and data export.
